% Data Protection Notice (DPDP).
%
% NOT derived from a .docx -- there was no data protection notice in the
% set. The Digital Personal Data Protection Act, 2023 and the DPDP Rules
% notified on 13 November 2025 require that a Data Principal be given
% notice of what is processed and why.
%
% Section 7 of the Act treats processing "for the purposes of employment"
% as a legitimate use, so the Intern's own data does not turn on consent.
% Emergency contacts are a different matter and are handled on the
% Emergency Contact Form. See DECISION-DEBT.md D-4.
%
% WORDING IS A DRAFT FOR REVIEW BY QUALIFIED COUNSEL.
\documentclass{paridyum}
\begin{document}

\paridyumletterhead
\paridyumdateref{\dkDocDate}{\dkDocRef}
\paridyumtitle{Data Protection Notice}

This notice explains what personal data \textbf{\dkOrgName} collects about
\textbf{\dkPIdentityLegalName} in connection with the internship, why it is collected, how
long it is kept, and what rights the Intern has.

\section*{1. What Paridyum Collects}
\begin{itemize}[leftmargin=*, noitemsep]
    \item identity and contact details --- name, preferred name, email address,
          telephone number and postal address;
    \item education details --- institution, programme, specialization, expected year
          of graduation and institutional email address;
    \item internship details --- role, department, supervisor, dates, working
          arrangement and weekly commitment;
    \item professional details voluntarily provided --- such as a portfolio,
          repository or professional profile;
    \item device and access details --- the device used for internship work, accounts
          issued, and access granted and removed;
    \item records of documents issued to and signed by the Intern; and
    \item emergency contact details, which are collected separately on the Emergency
          Contact Form.
\end{itemize}

Paridyum does not require and does not seek health records, financial account details,
government identity numbers, or biometric data for this internship.

\section*{2. Why Paridyum Processes It}
Personal data is processed for purposes connected with the internship, namely:
\begin{itemize}[leftmargin=*, noitemsep]
    \item arranging and administering the internship;
    \item issuing and managing system access, and keeping systems and information
          secure;
    \item preparing the documents that form the internship record, including
          agreements, policies, certificates and reference letters;
    \item providing documentation required by the Intern's educational institution;
          and
    \item meeting Paridyum's legal and record-keeping obligations.
\end{itemize}

Personal data is not used for marketing, is not sold, and is not used to train
external systems.

\section*{3. Who It Is Shared With}
Personal data is available to the Intern's supervisor and to those at Paridyum who
need it to administer the internship. It is shared outside Paridyum only where the
Intern asks for it --- for example, documentation sent to an educational institution
--- or where disclosure is required by law.

\section*{4. How Long It Is Kept}
The internship record is kept for the duration of the internship and for the retention
period that follows, as set out in Paridyum's retention schedule. Documents that
evidence what was agreed and issued --- the signed agreements and the certificate
record --- are kept for longer, because their purpose is to remain provable. Personal
data no longer needed for any of these purposes is deleted.

\section*{5. The Intern's Rights}
The Intern may:
\begin{itemize}[leftmargin=*, noitemsep]
    \item ask what personal data Paridyum holds about them and how it is processed;
    \item ask for inaccurate or incomplete data to be corrected or completed;
    \item ask for personal data to be erased where it is no longer required for the
          purposes above and Paridyum is not required to keep it; and
    \item raise a grievance about how their personal data has been handled.
\end{itemize}

\section*{6. How to Exercise Them}
Write to \textbf{\dkOrgEmail}. Paridyum will respond within a reasonable period. If the
Intern is not satisfied with the response, they may complain to the Data Protection
Board of India.

\section*{7. Keeping Data Accurate}
The Intern should tell Paridyum if any of the information above changes during the
internship, so that the record stays accurate.

\section*{8. Acknowledgement}
By signing below, the Intern confirms that they have received and read this notice.

\vspace{1.5em}
\paridyumofficers{\dkSigCeoName}{\dkSigCeoDesignation}{\dkSigCtoName}{\dkSigCtoDesignation}

\vspace{1.5em}
\paridyumcountersign{Intern}{\dkPIdentityLegalName}

\end{document}
